%======================================================================
\section{Introduction}
\label{sec:intro}
%======================================================================

LLM agents can now debate, critique, vote, and aggregate their outputs.  Yet an aggregate is only as useful as the behavior that generated it.  When producing a complete, checkable report is more expensive than a short answer or imitation, a multi-agent system needs a way to make that trade-off observable and consequential.  Accuracy alone cannot answer this question: it is a terminal outcome, whereas mechanism design concerns the choices made before the outcome is known.

We study this problem through \textbf{AEIC-MD} (AEIC-mediated mechanism design for deliberation).  AEIC structures each report as an \emph{Assumption}, \emph{Evidence}, \emph{Inference}, and \emph{Conclusion}.  The structure is not merely a prompting format: it is the observable report space on which a leader can announce a credit and adoption rule.  Having observed this public rule, followers choose low-effort reporting, complete AEIC reporting, verification, or imitation.  This timing gives a concrete, testable meaning to a leader--follower mechanism: the intervention is successful only if it changes logged follower actions, rather than assigning a desirable score after generation.

The project originally included topology selection and offline quality weighting.  We retain them as contextual baselines, but make a different causal claim here.  \emph{Topology controls information flow; AEIC-MD controls the incentives under which information is produced.}  The latter is the main line because it is what the frozen experiments directly identify.

Our formal comparison uses 450 paired questions from StrategyQA, MMLU-Hard, and TruthfulQA.  Relative to no mechanism, Static AEIC raises structured high-effort action from 0.46\% to 56.93\% and decision quality from 68.00\% to 82.22\% (+14.22pp, Holm-adjusted McNemar $p=1.23\times10^{-11}$).  Outcome-only incentives do not improve quality (67.33\%, $p=0.780$).  Thus, the evidence is not simply that ``rewards help'': rewards tied to observable, locally checkable reports alter behavior in a way that generic outcome rewards do not.

We then test the interpretation rather than only the headline number.  Reward and cost sweeps show monotone within-domain effort responses.  Manipulated evidence, social-only information, copy pressure, and counterevidence conditions test whether AEIC exposes behavior under distinct information environments.  Further experiments cover signal and history ablations, two additional random seeds, and three-to-five-agent panels.  Finally, a held-out GLM backend does \emph{not} reproduce the main quality gain, an essential limitation that locates the next problem---contract calibration across model families.

Our contributions are:
\begin{enumerate}
  \item A revealed-rule, logged-action formulation of AEIC-mediated multi-agent deliberation, distinguishing follower response from retrospective response scoring.
  \item Paired formal evidence that a Static-AEIC rule raises both high-effort choices and decision quality, while outcome-only incentives do not.
  \item Boundary and robustness evidence---including reward/cost response, controlled information conditions, seeds, panel sizes, ablations, and a transparently negative held-out-backend result---that defines the scope of the claim.
\end{enumerate}
